\section{The quantitative mass window and its prime-counting boundary}
\label{sec:quantitative-windows}

\status{QW1--QW9 prove the literal conditional \texttt{window\_lower}
from \texttt{SiegelMassRatioCM}, including the Euler factors, both local
indices, the Vandermonde comparison and the uniform passage to the actual
host. QW10--QW13 prove the weighted-to-unweighted norm-prime reduction
and the two scale consequences of the explicit short-interval input.
QW14--QW15 prove classical nonvanishing at one and the finite-conductor
real-zero margin. The global mass-ratio theorem, the compatible arithmetic
\texttt{LFam}, and its uniform weighted short-interval lower bound remain
separate inputs. No arbitrary short interval is supplied by qualitative
prime existence or by the finite-conductor margin.}

In this section $B=e^L$, $L=x^\epsilon$, $a=xB$, and $r=k+m-1$.
The letter $B$ here is the source scale \texttt{s.b}; in QW10--QW13
the distinct symbol $\mathcal B=4dM^2$ denotes the analytic conductor
bound. The finite host has $[E:\mathbb Q]=2f$. All genus masses, local
indices and Harish--Chandra coordinates are the literal ones on the
same supplied host and preliminary datum.

\paragraph{QW1: absolute logarithmic control of the actual Euler product.}
Let $E$ satisfy the source \texttt{IsCMUnram} condition and put
$f=[E:\mathbb Q]/2$. For a maximal ideal $\mathfrak p$ above a rational
prime $p$, write $N\mathfrak p=p^{d_{\mathfrak p}}$. The actual factor is
\[
 E_{\mathfrak p}(M)=
 \begin{cases}
 (1-(-1)^M(N\mathfrak p)^{-M/2})^{-1},&
       \overline{\mathfrak p}=\mathfrak p,\\
 (1-(N\mathfrak p)^{-M})^{-1/2},&
       \overline{\mathfrak p}\ne\mathfrak p.
 \end{cases}                                                   \tag{QW.1}
\]
If $\overline{\mathfrak p}=\mathfrak p$, the unramified CM condition
provides an integer $z$ with $z-\bar z\notin\mathfrak p$. Thus residue
conjugation is nontrivial. It would be the identity on a residue field
$\mathbb F_p$, so $d_{\mathfrak p}\ge2$. In all cases, using
$\log(1+t)\le-\log(1-t)$ for $0\le t<1$,
\[
 E_{\mathfrak p}(M)>0,\qquad
 |\log E_{\mathfrak p}(M)|
 \le\frac{d_{\mathfrak p}}2[-\log(1-p^{-M})].                    \tag{QW.2}
\]
Indeed the fixed-prime factor has absolute logarithm at most
$-\log(1-p^{-M})$; the other has at most half this bound.
The ideal factorization identity
$\sum_{\mathfrak p\mid p}e_{\mathfrak p}d_{\mathfrak p}=2f$ gives
$\sum_{\mathfrak p\mid p}d_{\mathfrak p}/2\le f$.
For $M\ge2$, unique factorization of positive integers gives
\[
 \sum_p-\log(1-p^{-M})=\log\zeta(M)<\infty.
\]
This last equality can be obtained by increasing finite Euler products:
their positive expansions enumerate the integers supported on those
primes, and their limit is $\sum_{n\ge1}n^{-M}$.
Therefore the sum of the absolute logarithms in (QW.2) converges.
Exponentiating its signed sum constructs the actual positive product
$L_E(M)=\prod_{\mathfrak p}E_{\mathfrak p}(M)$ and proves
\[
               |\log L_E(M)|\le f\log\zeta(M)\le f.             \tag{QW.3}
\]
For the last bound, the integral test gives
$\zeta(M)\le\zeta(2)\le2$, and $\log2\le1$.
This establishes the source Euler-product existence and logarithmic
bound; it makes no assumption about zeros of an analytically continued
$L$-function.

\paragraph{QW2: count the two local level groups before taking their ratio.}
For a finite field of cardinality $q\ge3$, choosing successive independent
columns gives
\[
 |\mathrm{GL}_n(\mathbb F_q)|=q^{n^2}\theta_q(n),\qquad
 \theta_q(n)=\prod_{j=1}^n(1-q^{-j}),\qquad
                 \tfrac12\le\theta_q(n)\le1.                   \tag{QW.4}
\]
The lower bound follows from
$\prod(1-t_j)\ge1-\sum t_j$ for $t_j\in[0,1]$ and
$\sum_{j\ge1}q^{-j}=1/(q-1)\le1/2$.
The column-level unitriangular group has $k(k-1)/2$ independent entries.
For the row level put $s=m/2-1$, with the actual even $m\ge2$ and
$r=k+m-1$. Its free entries comprise the $k$-by-$k$ unitriangular
part and the two outer blocks, giving exactly
\[
       s(2k+1+2s)+\frac{k(k-1)}2
       =\frac{r^2-(k+1)^2-2s}2+\frac{k(k-1)}2.                  \tag{QW.5}
\]
Every choice of these entries gives an upper triangular matrix with
diagonal one. Conversely the defining zero pattern leaves precisely
these entries free. Its cardinality is therefore the corresponding
power of $q$; no quotient by an additional determinant character occurs.
Dividing (QW.4) by these two subgroup cardinalities proves
\[
 \tfrac12q^{k(k+1)/2}\le I_{\rm col}(q)\le q^{k(k+1)/2},\qquad
 \tfrac12q^{(r^2+3k+1+2s)/2}\le I_{\rm row}(q)
                   \le q^{(r^2+3k+1+2s)/2}.                   \tag{QW.6}
\]
The exponents are integers for the actual parities. Their difference
is $m(2k+m-1)/2$. At the two host primes $g_0,g_1$, each with $f$
copies, (QW.6) gives positive indices and
\[
 \frac{I_{\rm row}}{I_{\rm col}}
 \ge2^{-2f}(g_0g_1)^{fm(2k+m-1)/2}.                            \tag{QW.7}
\]
This is the original \texttt{idx\_place}/\texttt{idx\_ratio} statement
for the actual level groups, including its factor $2^{-2f}$.

\paragraph{QW3: telescope the stated mass hypothesis without strengthening it.}
Write $\mu_N$ for the genus mass of the standard rank-$N$ lattice.
IG9 proves its positivity directly from the finite positive mass sum.
The retained \texttt{SiegelMassRatioCM} hypothesis says, for every
unramified CM field and every $N\ge3$,
\[
 \frac{\mu_N}{\mu_{N-1}}
 =R_N:=\frac{\big((N-1)!(v/\pi)^N\big)^f}{\sqrt{(2v)^f}}L_E(N),
 \qquad v=\frac{|\operatorname{disc}E|^{1/(2f)}}2.               \tag{QW.8}
\]
For $k\ge2$ and $r=k+m-1$, multiply these identities, cancelling the
positive intermediate masses, to obtain
$\mu_r=\mu_k\prod_{N=k+1}^rR_N$. If $m=1$ the product is empty and
the identity is $\mu_k=\mu_k$. In the actual host, the identities
$f\mathrm{Deg}(E)=f$ and $v\mathrm{Inv}(E)=v$ follow from
$[E:\mathbb Q]=2f$ and total complexity, so (QW.8) has exactly the
normalization required by \texttt{mass\_tel}.

This paragraph proves the telescope from the displayed ratio law. The
ratio law itself is not a consequence of positivity. Its pinned
producer passes through \texttt{HermRepUnram} and \texttt{LIStat},
the representation and local-density mass identities. Those are
retained source targets.

\paragraph{QW4: an elementary Vandermonde comparison for interlacing columns.}
For a decreasing $N$-tuple $z$ put
$V_N(z)=\prod_{0\le i<j<N}(z_i-z_j)$ and
$d_N(z)=V_N(z)/\prod_{i=0}^{N-1}i!$.
Suppose $k\ge2$, $v_i\ge0$, $v_i-v_{i+1}=4\Delta$ with
$\Delta\ge1$, and
$u_i>v_i>u_{i+1}$. Let $b_0<\cdots<b_{m-1}$ be positive with
successive gaps at least one, and append $-b_0,\ldots,-b_{m-1}$ to $v$.
Every cross factor $v_i+b_j$ is at least $b_j$ and every bottom factor
$b_j-b_i$ is at least one. Hence
\[
 V_r(v,-b)\ge V_{k-1}(v)\prod_{j=0}^{m-1}b_j^{k-1}.             \tag{QW.9}
\]
If every gap in $u$ is at most $D\ge1$, the $2k-3$ pairs involving
its first or last entry contribute at most $D^{2k-3}$. Every remaining
pair satisfies
$u_i-u_j\le v_{i-1}-v_j$ for $1\le i<j\le k-2$.
These are exactly the nonadjacent pairs of $v$. The omitted adjacent
factors of $V_{k-1}(v)$ are at least one, so
\[
                      V_k(u)\le D^{2k-3}V_{k-1}(v).             \tag{QW.10}
\]
The proof covers $k=2$ by empty products. Thus this comparison loses
only $2k-3$ endpoint factors, rather than the full number of pairs.

For the literal \texttt{PlaceData}, $1\le x$, $a=xB$, and
$v_i\le4k\Delta$. Either $u_0\le2a$, $u_{k-1}\ge-a$, or
$u_0\le12k\Delta$, $u_{k-1}=v_{k-2}-1/2\ge-1/2$.
The former gives total gap at most $3a$. In the latter use
$4k\Delta\le b_0/x$ and $b_0<w_0\le w_{m-1}\le2B$ to bound it by
$6B+1/2\le8a$. Taking $D=8a$ in (QW.9)--(QW.10) proves, in particular,
\[
 \log V_r(v,-b)-\log V_k(u)
 \ge(k-1)\sum_j\log b_j-2k\log(8a)-2k(1+\log k).              \tag{QW.11}
\]
All Vandermondes are positive. The last extra nonpositive term merely
weakens our bound to the form used by the pinned source.

\paragraph{QW5: pass to the middle weights with an explicit constant.}
The remaining \texttt{PlaceData} clauses say
$b_j<w_j<b_{j+1}$ where defined, $w_{m-1}\le2B$, and
$4k\Delta\le b_0/x$. Thus $b_0\ge1$ and
\[
 \sum_j\log(w_j/b_j)
 \le\log(w_{m-1}/b_0)\le1+L,\qquad
                     \sum_j\log b_j\le m(1+L).               \tag{QW.12}
\]
The first inequality telescopes, including the case $m=1$.
For $\epsilon>0$, $L=x^\epsilon\ge1$ and $k\le x^{C_0}$,
\[
 \log x\le L/\epsilon,\quad
 \log k\le C_0L/\epsilon,\quad
 \log(8a)\le7+L/\epsilon+L.
\]
Insert (QW.12) into (QW.11). The loss is at most
\[
 k(1+L)+m(1+L)+2k(7+L/\epsilon+L)
                 +2k(1+C_0L/\epsilon).
\]
Since $L\ge1$, this proves the exact useful one-place bound
\[
 k\sum_j\log w_j\le\log V_r(v,-b)-\log V_k(u)
       +kL\left(20+\frac{2+2C_0}{\epsilon}\right)+2mL.          \tag{QW.13}
\]
No asymptotic estimate for individual gamma factors is hidden here.

\paragraph{QW6: factorial cancellation gives the global logarithmic ratio.}
Suppose $f,m\ge1$, $k\ge2$, and every one of the $f$ places satisfies
the same \texttt{PlaceData}. Assume $v\ge\pi$, $\log v\le L$,
$g_0,g_1\ge1$, the positive masses and indices, (QW.3), (QW.7),
and the mass telescope (QW.8). Retain the exact hypotheses
\[
 d_* =\mu_k I_{\rm col}\prod_\rho d_k(u_\rho),\qquad
 D_*\ge(2/3)^f\mu_r I_{\rm row}\prod_\rho d_r(v_\rho,-b_\rho).
                                                                    \tag{QW.14}
\]
These already give $d_*,D_*>0$. In the logarithm of the ratio, the
factorial terms cancel exactly, because
\[
 \sum_{N=k+1}^r\log((N-1)!)+\sum_{i=0}^{k-1}\log(i!)
                         =\sum_{i=0}^{r-1}\log(i!).            \tag{QW.15}
\]
Put $\ell=\log(v/\pi)\in[0,L]$, $G=\log g_0+\log g_1$,
$T=\sum_{\rho,j}\log w_{\rho j}$, and $K=(2+2C_0)/\epsilon$.
The remaining mass contribution is bounded below by
\[
 fk(m-1)\ell-\tfrac f2(m-1)(1+L)-f(m-1),
\]
since $\sum_{N=k+1}^rN\ge k(m-1)$, $\log2\le1$, and (QW.3).
The index contribution is at least $fkmG-2f\log2$; the Vandermonde
contribution is at least $kT-fkL(20+K)-2fmL$ by (QW.13).
Finally $f\log(2/3)\ge-f/2$.
Relative to the desired main term $kT+fkm\ell+fkmG$, the total loss
divided by $f$ is no more than
\[
 kL(20+K)+2mL+k\ell+\tfrac{m-1}{2}(1+L)
                     +(m-1)+2\log2+\tfrac12.
\]
Using $L\ge1$, $m\ge1$, $k\ge2$, this is at most
$(30+K)(kL+m^2L)$. Consequently the literal
\texttt{LogRatioLowerBound} holds with the explicit choice
\[
 \boxed{C_\epsilon=30+(2+2C_0)/\epsilon},                       \tag{QW.16}
\]
and conclusion
\[
 \frac{\log(D_*/d_*)}{fk}\ge
 \frac1f\sum_{\rho,j}\log(w_{\rho j}v/\pi)
 +m(\log g_0+\log g_1)-C_\epsilon\left(L+\frac{m^2L}{k}\right).
                                                                    \tag{QW.17}
\]
The restrictions $0<\epsilon<1/52$ and $C_0\ge0$ of the source are
preserved, although this calculation uses only $\epsilon>0$ and the
displayed size assumptions.

\paragraph{QW7: verify PlaceData on the actual top, thick and thin places.}
The scale and rank estimates RF11--RF14 were proved with thresholds
independent of the eventual datum. Fix $c_D,c_h\ge0$ and set
$\Delta=\lceil x^{c_D}\rceil$, $H_0=2\lceil x^{c_h}\rceil$.
They give, uniformly in the actual rank record,
\[
 x(4k\Delta+H_0)\le w''\le B',\quad 3B'\le B,\quad
 mH_0+8f^2+2\le B,\quad k\le x^5.                             \tag{QW.18}
\]
For completeness, $\Delta\le2x^{c_D}$, $H_0\le4x^{c_h}$,
$k\le x^5$, $f\le x^3$, $m\le x$ bound each polynomial on the left.
RF11 makes $w''$ and $B$ exceed any prescribed such polynomial,
uniformly for $1\le n\le x^3$, and gives $B'/B=x/X'\to0$.
The actual bottom-window length is at most
$2n x^{-\epsilon/8}L\Lambda B'/k\le B'$. Indeed RF14 gives
$\Lambda\le C L(\log x)^2$. Use also the sharper original rank
inequalities $n\le2x^{2+8\epsilon}$ and
$k\ge x^{9/2+10\epsilon}$; the coefficient is then at most
$C' x^{-5/2-\epsilon/8}(\log x)^2$, which tends to zero.

Write the actual bottom start at a place as $\beta$. The \texttt{HCVal}
equalities are
\[
 v_i=(k-2-i)4\Delta,\qquad b_j=\beta+jH_0,\qquad
                         w_j=b_j+H_0/2.                       \tag{QW.19}
\]
At a top place, $\beta=b_T+\lfloor B'/x\rfloor$ and the original
window plus the preceding length bound puts $B'\le b_T\le2B'$.
At a thick place, $\lceil B\rceil\le\beta<\lceil B\rceil+
|\mu(E)|$. At a thin place, $\beta=\lceil w''\rceil$.
The source root-of-unity bound is
$|\mu(E)|\le2[E:\mathbb Q]^2=8f^2$: if its order is $q$, then
$\varphi(q)\le[E:\mathbb Q]$ and $q\le2\varphi(q)^2$.
The latter follows from the prime-power formula for $\varphi$:
every prime-power contribution to $\varphi(q)^2/q$ is at least one
except a single possible factor $1/2$ from $2^1$.
These three cases therefore give
\[
                   w''\le\beta\le B+8f^2+2.                  \tag{QW.20}
\]
Together (QW.18)--(QW.20) imply
$x(4k\Delta+H_0)\le\beta$ and $\beta+mH_0\le2B$.

The existing \texttt{HCParam} gaps give the strict interlacing of $u,v$,
$-b_0<u_{k-1}$, $b_j<w_j<b_{j+1}$ and $b_{j+1}-b_j\ge1$.
Equation (QW.19) gives $|w_j-b_j|=H_0/2\le b_j/x$,
$4k\Delta\le b_0/x$, $w_{m-1}\le2B$, and $w_j\ge1$.
At a top place the prescribed endpoints give $u_0\le2a$ and
$u_{k-1}\ge-a$; the small endpoint correction is at most
$a/(k x^{A/2})\le a/2$. At another place the exact endpoints are
$u_0=v_0+4k\Delta+1/2$, $u_{k-1}=-1/2=v_{k-2}-1/2$.
Since $4k\Delta\ge1$, they meet the alternative endpoint inequalities.
This checks all nine original \texttt{PlaceData} clauses and $w_j\ge1$
on the same datum, at every word and every place.

\paragraph{QW8: the host constants and rank losses are uniformly small.}
The actual host discriminant inequalities give
\[
 \operatorname{rd}(E)\ge
 2^{e+1-2^{-e}}\sqrt d\ge4\sqrt3\ge32/5,
 \qquad v=\operatorname{rd}(E)/2\ge16/5>\pi,
\]
using $e\ge2$ and $d\ge3$. The other host inequality and its supplied
prime-window bounds give
\[
 \operatorname{rd}(E)\le(Q_S/2)\sqrt{q d}\le Q_S qd,
 \quad \log(Q_S qd)\le2x^{\epsilon/5}(\log x)^2\le L
\]
eventually. The last step is RF11 with positive exponent gap
$4\epsilon/5$. Hence $\log v\le L$.
Also the same \texttt{HostVal} has $g_i\ge(kB)^8$, so
$\log g_i\ge8L$ and $g_i\ge3$ for large $x$.

For any fixed real $C,C_\Gamma$, RF12 gives, eventually,
\[
      2|C|\le m,\qquad m^2\le k,\qquad 2C_\Gamma m^4\le k,
      \qquad L\ge1.                                           \tag{QW.21}
\]
One explicit way to see the margins is
$m\ge x^{1/4}$, $m\le x$, $k\ge x^{9/2+10\epsilon}$.
After $2|C_\Gamma|\le x^{1/4}$, the third left side is at most
$x^{17/4}\le k$; the other inequalities follow at once.
All constants were fixed before this final threshold. Neither the word,
the host, nor the chosen coordinates affect it.

\paragraph{QW9: prove the original quantitative window with its quantifier order.}
Fix \texttt{SiegelMassRatioCM}, $0<\epsilon<1/52$, a standard constant
record $\mathcal C$, and $c_D\ge0$. Choose $c_2=0$. For each
$c_h\ge c_2$, take a common $x_0$ for QW7--QW8, with
$C=C_\epsilon=30+12/\epsilon$ in (QW.21). Now let any scale $s$,
any original \texttt{PreCh} $P$, and the remaining $q_1,a,b,c,du,D_s$
be supplied with $s.\epsilon=\epsilon$, $x\ge x_0$, the literal
\texttt{ChOK}, and the two exact size hypotheses (QW.14).
For each word, QW1--QW8 apply with $C_0=5$.
The rows' $v,b$ coordinates agree with those of the base word by the
original \texttt{hcW\_v}/\texttt{hcW\_b} identities; thus the single
given lower bound for $D_s$ is exactly the one needed for every word.
Enumerating $\mathrm{Tau}$ replaces its product by the $f$ place product
without changing any factor.

The average logarithm in (QW.17) is nonnegative, since $w_j\ge1$ and
$v\ge\pi$. Moreover
\[
 m(\log g_0+\log g_1)\ge16mL,\qquad
 C_\epsilon(L+m^2L/k)\le2C_\epsilon L\le mL.
\]
Consequently $\log(D_s/du)/(fk)\ge15mL$. Put
$\Gamma=\exp(C_\Gamma m^5fL)$, the exact source \texttt{GamS}.
Positivity from (QW.14) justifies the real powers and logarithms, and
\[
 \begin{split}
 \log\!\left((D_s/du)^{1/(fk)}\Gamma^{-2/(fk)}\right)
 &=\frac{\log(D_s/du)}{fk}-\frac{2C_\Gamma m^4}{k}mL\\
 &\ge14mL\ge mL.
 \end{split}                                                    \tag{QW.22}
\]
This is \texttt{window\_lower}, with some slack preserved. The threshold
was chosen before $P,q_1,a,b,c,du,D_s$ and the word. This proves the
literal conditional theorem, not the existence of all its compatible
geometric inputs, and not the retained mass-ratio law itself.

\paragraph{QW10: remove the prime powers from the actual weighted sum.}
For the source family $\mathcal F=\texttt{LFam}$, retain its coefficient
laws $\lambda_g(n)\ge0$, $\sum_g\lambda_g(n)\le2\Lambda(n)$, and
the arithmetic support statement
\[
 p\text{ prime},\ \lambda_0(p)\ne0
                  \quad\Longrightarrow\quad\mathrm{NormOne}(d,M,p).
                                                                    \tag{QW.23}
\]
The latter is a condition on the same family, not merely an unrelated
prime represented by a norm. Split the sum over integers in $(X,Y]$
into primes and nonprimes. A prime outside the norm-prime set has
$\lambda_0(p)=0$ by (QW.23), whereas a prime inside it has
$\lambda_0(p)\le2\log p\le2\log Y$. At nonprimes the bound by
$2\Lambda(n)$ leaves only prime powers of exponent at least two.
For $Y\ge1$ this proves exactly
\[
 \sum_{\lfloor X\rfloor<n\le\lfloor Y\rfloor}\lambda_0(n)
 \le2\log Y\,N_{d,M}(X,Y)+2(\psi(Y)-\vartheta(Y)),               \tag{QW.24}
\]
where $N_{d,M}$ is the cardinality of the original norm-prime set.
The nonnegative remainder was enlarged to the full interval $(0,Y]$.
Elementary Chebyshev estimates give
$\psi(Y)-\vartheta(Y)\le C\sqrt Y$ for one absolute $C$.
For example the binomial coefficient estimate first gives
$\vartheta(t)\le C_1t$; then
$\sum_{j\ge2}\vartheta(Y^{1/j})\le C_1\sqrt Y+
C_1(\log_2Y)Y^{1/3}\le C_2\sqrt Y$ after enlarging the constant.
This uses no prime-number-theorem error term.

\paragraph{QW11: isolate exactly the analytic input that supplies ShortNorm.}
Retain the source arithmetic construction of an $\mathcal F$ satisfying
(QW.23) with conductor bound $\mathcal B=4dM^2$, for
$d$ prime, $d\equiv3\pmod4$, $d\ge100$, $M\ge3$.
Retain also its uniform weighted estimate: there are $A,c_0,e>0$ such
that, for every such family and $X\ge1$,
\[
 \log\mathcal B\le c_0\sqrt{\log X},\quad
 e^{-c_0\sqrt{\log X}}\le h\le1
 \quad\Longrightarrow\quad
 \sum_{X<n\le X(1+h)}\lambda_0(n)\ge\frac{e hX}{\mathcal B^A}.
                                                                    \tag{QW.25}
\]
Here and below integer endpoints mean the same floor convention as
(QW.24). Set $K=8(|C|+1)/e$ and choose
\[
 c=\min\left(c_0,\frac1{2(1+A)+2+K}\right)>0.
\]
For $s=\sqrt{\log X}$, $\mathcal B\ge4$ implies
$1\le\log\mathcal B\le cs$, so $s\ge2(1+A)+2+K$.
Consequently $c\le1$, $\mathcal B^A\le e^{As}$,
$h\ge e^{-s}$, and
\[
 K\le e^{s^2/2-(1+A)s},\qquad
                       K\sqrt X\le hX/\mathcal B^A.           \tag{QW.26}
\]
The first inequality follows from $e^t\ge1+t$ and the displayed
lower bound for $s$; multiplying the three exponential bounds proves
the second. In particular $\log X\ge1$.

For $Y=X(1+h)\le2X$, one has
$\log Y\le2\log X$ and $\sqrt Y\le2\sqrt X$.
Equations (QW.24)--(QW.26) show that the prime-power remainder consumes
at most half the right side of (QW.25). Dividing by $4\log X$ gives
\[
       N_{d,M}(X,X(1+h))\ge
          \frac{(e/8)hX}{\mathcal B^A\log X}.                  \tag{QW.27}
\]
This reconstructs the complete reduction to \texttt{ShortNorm}, with
$c'=e/8$. The arithmetic \texttt{ExistsLFamRayStatement} and analytic
\texttt{TrivialClassShortIntervalStatement} suppliers of (QW.23)--(QW.25)
remain; their names do not replace proofs of those assertions.

\paragraph{QW12: derive the source level-prime window uniformly in its endpoint.}
Assume \texttt{ShortNorm}, namely (QW.27) with some fixed $A,c,c'>0$
under its displayed hypotheses. The original \texttt{ModAt} says
\[
 d\text{ prime},\ d\equiv3\pmod4,\ d\ge100,\quad
 N>0,\ d\mid N,\quad \log N\le3x^{\epsilon/5}(\log x)^2.       \tag{QW.28}
\]
Fix $0<\epsilon<1/52$, $c_s,C_Y>0$. Suppose further
$g_i\ge1$, $\log g_i\le C_Yx^\epsilon$, $m\ge x^{1/3}$,
$1\le V\le x^{12}$, and $T_s>0$ with
$\log T_s\ge c_s m x^\epsilon$.
Set $M=N(g_0g_1)^2$ and $h=1/V$. Then $M\ge N\ge d\ge100$, and
\[
 \log(4dM^2)\le\log4+3\log N+8C_Yx^\epsilon,
 \qquad \sqrt{\log T_s}\ge\sqrt{c_s}x^{1/6+\epsilon/2}.
\]
Both $x^{\epsilon/5}(\log x)^2$ and $x^\epsilon$ are little-oh of
$x^{1/6+\epsilon/2}$; the second uses $\epsilon<1/3$.
Also $\log V\le12\log x$ is little-oh of that same power.
RF11 therefore supplies a single threshold, depending only on the fixed
constants, for which
$\log(4dM^2)\le c\sqrt{\log T_s}$ and
$h\ge e^{-c\sqrt{\log T_s}}$.
The positive right side of (QW.27) makes the finite norm-prime set
nonempty. Every chosen member satisfies exactly
\[
 T_s<T\le(1+1/V)T_s,\qquad
 T\text{ prime},\quad \mathrm{NormOne}(d,N(g_0g_1)^2,T).        \tag{QW.29}
\]
The floor endpoints give the strict lower inequality in (QW.29).
This proves \texttt{levelPrimeN} with the source quantifier order:
$T_s,V,m,N,d,g_i$ are chosen after the common threshold. Arbitrarily
large $T_s$ only improve the two tested inequalities. For the actual
host one may insert $C_Y=9$, $V=fk\le x^{12}$ and the original endpoint
lower bound to obtain the corresponding conditional \texttt{LevelPrime}.
The unconditional producer still needs QW11's retained suppliers.

\paragraph{QW13: the same input gives the literal nine-tenths dyadic count.}
Assume \texttt{ShortNorm} and \texttt{ModAt} again, now with any fixed
$\epsilon>0$. Put $t=x^{\epsilon/5}$, $\mathcal B=4dN^2$ and
$y=\log Y\ge8x^\epsilon=8t^5$. For sufficiently large $x$,
\[
 \log\mathcal B\le\log4+3\log N
 \le\log4+9t(\log x)^2\le10t^2.                               \tag{QW.30}
\]
Thus $\log\mathcal B/\sqrt y\le(10/\sqrt8)t^{-1/2}\le c$.
Apply (QW.27) with $h=1$ to obtain
$N_{d,N}(Y,2Y)\ge c'Y/(\mathcal B^A y)$.
The inequality needed for this to exceed $Y^{9/10}$ is
\[
                A\log\mathcal B+\log y-\log c'\le y/10.       \tag{QW.31}
\]
For $y\ge10$, $y/10-\log y$ is increasing. Hence it is enough to
check $y=8t^5$. At that value (QW.30) reduces (QW.31) to
$10At^2+\log(8t^5)-\log c'\le(4/5)t^5$, true eventually.
This proves, uniformly over every $Y$ with the original lower bound,
\[
                         N_{d,N}(Y,2Y)\ge Y^{9/10}.             \tag{QW.32}
\]
The exponent $9/10$ and the condition $\log Y\ge8x^\epsilon$ are
the literal \texttt{normPrimesN} output. They are conditional on the
quantitative analytic input, not consequences of IG's qualitative
unboundedness of ray primes.

\paragraph{QW14: a nonnegative character product proves nonvanishing at one.}
For completeness we supply the classical nonvanishing used by the
finite-modulus source lemma. Fix a modulus $N\ge1$. Use the standard
analytic continuation of Dirichlet $L$-functions: a nonprincipal
$L(s,\chi)$ is entire, and the principal factor is
$\zeta(s)\prod_{p\mid N}(1-p^{-s})$, with a single simple pole at one
and a zero at $s=-2$. These are the classical continuation and trivial
zero facts, independent of RH.\footnote{For the stated continuation,
principal pole and trivial-zero formulas, see NIST DLMF,
\url{https://dlmf.nist.gov/25.15}. The argument below supplies the
nonvanishing endpoint used by the pinned source's \texttt{noExc\_one}.}
Suppose some nonprincipal factor vanished at one. Then
\[
                  F(s)=\prod_{\chi\bmod N}L(s,\chi)
\]
would be entire, and $F(-2)=0$. In $\Re s>1$ its Dirichlet coefficients
are nonnegative: at a prime $p\nmid N$, if its class has order $d_p$
in $G=(\mathbb Z/N\mathbb Z)^\times$, character orthogonality gives
\[
 \prod_\chi(1-\chi(p)T)^{-1}=(1-T^{d_p})^{-|G|/d_p}.
\]
Every coefficient is nonnegative and the constant coefficient is one.
Write $F(s)=\sum_{n\ge1}a_n n^{-s}$, $a_n\ge0$, $a_1=1$.
For a fixed real $\sigma>1$, termwise differentiation is justified by
absolute convergence in a smaller right half-plane, and
\[
 (-1)^jF^{(j)}(\sigma)=\sum_{n\ge1}a_n n^{-\sigma}(\log n)^j\ge0.
\]
The Taylor series of an entire function has infinite radius. Evaluating
at $\sigma-t$ with $t=\sigma+2>0$ gives a sum of nonnegative terms,
whose constant term is $F(\sigma)\ge1$. Thus $F(-2)\ge1$, contradicting
$F(-2)=0$. Therefore every nonprincipal $L(1,\chi)$ is nonzero.
The argument establishes this classical endpoint without any assertion
about nontrivial zeros elsewhere in the critical strip.

\paragraph{QW15: finite-conductor real-zero margins preserve the final quantifier.}
For each fixed nonprincipal $\chi$, continuity and QW14 provide
$c_\chi>0$ such that $L(\beta,\chi)\ne0$ whenever
$1-c_\chi<\beta<1$ is real. Assign the harmless value $c_\chi=1$
to the principal character, for which no assertion is required.
There are finitely many characters modulo each positive $N$, and
finitely many $N\le N_0$. Their minimum is positive. It follows that
\[
 \forall N_0\ \exists c_0>0\ \forall\,1\le N\le N_0\
 \forall\chi\ne1\ \forall\beta\in(1-c_0,1):\quad
                              L(\beta,\chi)\ne0.               \tag{QW.33}
\]
For $N_0=0$ the assertion is vacuous and $c_0=1$ works. This is exactly
\texttt{QH.small\_zeros} and its literal real-variable \texttt{NoExc}.
The constant can depend on $N_0$; the proof provides no effective rate,
no fixed constant for all conductors, and no complex zero-free box.
In particular it cannot replace the uniform estimate (QW.25), the
independent simultaneous-prime selection, or the final RH/GRH argument.
